\section{随机变量之和与卷积}
\label{sec:05-Probability/05-Joint-Distributions/03}

两颗骰子。你想知道点数之和等于 7 的概率。不可能只盯着联合表格的某一个格子——和为 7 来自 \((1,6),(2,5),\ldots,(6,1)\)，六条完全不同的路径。你需要沿着直线 \(x+y=7\) 在联合表格里扫过去，把所有符合条件的格子概率加在一起。

这就是随机变量求和的核心问题：和的分布不是去联合分布的某个单点看一眼，而是要汇总一整个方向上的所有可能分解。

\subsection{离散卷积：沿对角线求和}

设 \(S=X+Y\)。对于 S 可能取到的每一个值 s，把所有满足 \(x+y=s\) 的联合概率加起来：

\[
P(S=s)
=\sum_x P(X=x,\,Y=s-x).
\]

这个公式不需要 X 和 Y 独立——只要你知道联合分布，就可以直接套。求和遍历 x，每取一个 x，y 就被锁死为 \(s-x\)。你遍历的是直线 \(x+y=s\) 上所有格点。

如果 X 和 Y 恰好独立，联合可以拆成边缘乘积：

\[
p_S(s)
=\sum_x p_X(x)\,p_Y(s-x).
\]

右侧这个求和形式叫离散卷积，记作

\[
p_S = p_X * p_Y.
\]

独立性的唯一作用就是把联合概率拆开。如果变量不独立，用回第一条带联合概率的公式。

拿两个独立 Bernoulli(\(p\)) 做例子。X 和 Y 各自以概率 p 取 1，\(1-p\) 取 0。和 S 可以取 0、1 或 2：

\[
\begin{aligned}
P(S=0) &= P(X=0,Y=0) = (1-p)^2,\\[2pt]
P(S=1) &= P(X=1,Y=0) + P(X=0,Y=1) = 2p(1-p),\\[2pt]
P(S=2) &= P(X=1,Y=1) = p^2.
\end{aligned}
\]

\(S=1\) 的系数 2 来自两条路径：X 成功且 Y 失败，以及 X 失败且 Y 成功。如果变量不独立，这两条路径的概率就不一定相等，公式会复杂一些，但求和的思想不变。

\subsection{连续卷积：沿直线积分}

连续情形怎么办？X 和 Y 有联合密度。固定一个和值 s，所有满足 \(x+y=s\) 的点构成平面上的一条斜线。你不能像离散那样一个个点求和——直线上的点有无穷多个，每个单点的概率都是零。

但沿着这条直线，联合密度有值。你的工作是：沿着 \(y=s-x\) 这条参数化直线，把联合密度对参数 x 积分。如果 X 和 Y 独立，

\[
f_S(s)
=\int_{-\infty}^{\infty}
f_X(x)\,f_Y(s-x)\,dx.
\]

几何上看，固定 s 后，你穿过联合密度曲面，沿着斜线 \(y=s-x\) 切出一道截面，然后把截面上各点的密度值沿 x 方向积起来。这不是弧长积分——参数是 x，积分元是 \(dx\)。精确的变换因子（Jacobian）会在\hyperref[sec:05-Probability/05-Joint-Distributions/05]{``多维变量变换与 Jacobian''}一节中作为变量变换加边缘化的结果重新推出。目前先接受这个公式作为离散卷积的连续对应物。

\subsection{两个均匀分布的和：三角形密度}

设 X 和 Y 独立，且各自在 \((0,1)\) 上均匀分布。联合密度在单位正方形上恒为 1，其余为零。求 \(S=X+Y\) 的密度。

关键在于：固定 s 之后，x 不能随便取。两个约束同时锁死：x 必须在 0 到 1 之间，且 \(s-x\) 也必须在 0 到 1 之间。这等价于 \(0<x<1\) 且 \(s-1<x<s\)。两个区间的交集才是真正有效的积分范围。不同的 s，交集形状不同。

对 \(0 < s < 1\)：交集为 \((0, s)\)。

\[
f_S(s)=\int_0^s 1\cdot 1\,dx = s.
\]

对 \(1 \le s < 2\)：交集为 \((s-1, 1)\)。

\[
f_S(s)=\int_{s-1}^1 1\cdot 1\,dx = 2-s.
\]

在其他范围（\(s \le 0\) 或 \(s \ge 2\)），两个支撑区间不相交，积分为零。结果是三角形密度：

\[
f_S(s)=
\begin{cases}
s, & 0<s<1,\\
2-s, & 1\le s<2,\\
0, & \text{其他}.
\end{cases}
\]

两个均匀独立变量加起来，密度从矩形变成了三角形。s=1 处的最大密度不是偶然——和为 1 的可分解方式最多（有整整一段 x 取值），和为 0.1 或 1.9 的分解路径短得多。密度的高度反映了``有多少种方式能达到这个和''。

如果你不看 s 的具体范围，直接写 \(\int_{-\infty}^\infty 1\cdot 1\, dx\)，积分会发散。积分上下限必须来自两个支撑区间的交集，而不是照抄 \(-\infty\) 和 \(\infty\)。这是连续卷积最容易出错的地方。

\subsection{差、线性组合与平均}

差 \(D = X - Y\) 可以理解成 \(X + (-Y)\)。如果 X 和 Y 独立且有密度，把卷积公式里的 \(y=s-x\) 换成 \(y=x-d\)（因为 \(d=x-y\) 即 \(y=x-d\)），就得到：

\[
f_D(d)
=\int_{-\infty}^{\infty} f_X(x)\,f_Y(x-d)\,dx.
\]

更一般的线性组合 \(aX+bY\) 可以通过变量变换、特征函数或者直接积分来求分布。但很多时候你不需要完整的分布——只需要前两阶矩：

\[
\E[X+Y]=\E[X]+\E[Y],
\]

\[
\Var(X+Y)
=\Var(X)+\Var(Y)
+2\Cov(X,Y).
\]

均值和方差公式只用到了各变量的矩和协方差，不需要做卷积。如果问题只问期望和方差，先去做完整卷积就是多余的工作。但在需要尾部概率、分位数或完整分布形状时，矩不够用，卷积绕不开。

\subsection{卷积为什么反复出现}

很多实际量是独立贡献之和：一天的总订单数、电路里的总噪声电压、一个请求在各服务节点的总延迟、累积降雨量。卷积就是概率论里的``加法规则'': 要求和的分布，就把各部分分布卷在一起。

有一类分布家族在卷积下具有封闭性——卷完之后还在同一个家族里，只是参数变了：

\begin{itemize}
\tightlist
\item
  独立 Poisson 变量之和仍是 Poisson，参数相加；
\item
  同尺度的独立 Gamma 变量之和仍是 Gamma，形状参数相加；
\item
  独立 Gaussian 变量之和仍是 Gaussian，均值和方差分别相加。
\end{itemize}

这些封闭性意味着，如果你知道各部分属于这些家族，就不需要每次重新做积分。但现实中的分布大多不封闭——卷两个不同家族的密度，通常只能硬算或数值求解。

另外，卷积严重依赖独立性假设。比如一次网络请求的总延迟是域名解析、传输、排队、数据库查询各阶段用时之和。如果各阶段近似独立，你可以逐次卷积它们的延迟分布。但在高负载场景，所有阶段往往同时受拥塞影响：独立性的失效意味着继续卷边缘分布会低估``所有阶段一起慢''的尾部风险。此时你需要先给定系统负载状态再建模，或者直接处理联合分布。

\hyperref[sec:05-Probability/07-Transforms-and-Bounds/01]{``生成函数、矩母函数与特征函数''}会提供另一个视角：傅里叶/拉普拉斯变换把卷积变普通乘法，好处是多项之和可以一口气处理。但变换方法和直接卷积积分的数学等价性没有改变——卷积仍然是底层机制。
